% ============================================================
% III. METODOLOGÍA
% ============================================================

\section{Metodología}

El procedimiento experimental se dividió en tres grupos: animales, instrumentos musicales y voces humanas. Cada archivo de audio se procesó con el paquete \texttt{animal\_analyzer} del proyecto \texttt{WORK\_THREE}, apoyado en las funciones genéricas de \texttt{signal\_lab/spectral.py} y en NumPy, SciPy \cite{scipy} y Matplotlib.

Las ocho grabaciones (tres animales, tres instrumentos, dos voces) se procesaron con \texttt{scripts/generate\_animal\_figures.py}, un script batch que llama directamente a \texttt{load\_audio} y \texttt{analyze\_audio} — las mismas funciones que usa la interfaz gráfica \texttt{run\_animal\_gui.py} para sus pestañas de animales, de modo que ambas vías producen resultados idénticos. Los animales y los instrumentos son audio real con licencia libre descargado de Wikimedia Commons; las dos voces son habla sintetizada localmente (Windows SAPI), usada como sustituto reproducible ante la falta de grabaciones humanas reales. La procedencia, autoría, licencia y limitaciones de cada archivo se documentan en \texttt{audio\_samples/SOURCES.md}.

El procedimiento general, tal como lo implementa \texttt{analyze\_audio}, fue:

\begin{enumerate}

\item Cargar el archivo de audio con \texttt{load\_audio}, que usa \texttt{librosa.load(path, sr=None, mono=True)} y preserva la frecuencia de muestreo nativa del archivo.

\item Registrar la frecuencia de muestreo y la duración total de la grabación.

\item Acotar la señal a los primeros $3$ s (\texttt{analysis\_window}, \texttt{ANALYSIS\_WINDOW\_S = 3.0}), dado que las grabaciones originales tienen una duración aproximada de $10$ s y solo esa ventana inicial se somete a análisis espectral y estadístico.

\item Representar la señal ventaneada en el dominio temporal.

\item Calcular la FFT de un solo lado con \texttt{compute\_fft} (\texttt{numpy.fft.rfft}, magnitud normalizada por $2/N$).

\item Obtener el espectro de magnitud y graficarlo.

\item Determinar la frecuencia dominante con \texttt{dominant\_frequency}, ignorando la componente en $0$ Hz.

\item Calcular la media y la desviación estándar con \texttt{signal\_stats}.

\item Calcular el puntaje de separabilidad estadística (Ec.~\ref{eq:separabilidad}) entre cada par de señales con \texttt{separability\_score} / \texttt{compare\_categories}.

\item Comparar los resultados obtenidos entre las diferentes señales.

\end{enumerate}

\subsection{Señales animales}

Se analizaron tres sonidos correspondientes a un gato, un perro y un gallo, cargados desde \texttt{audio\_samples/gato/}, \texttt{audio\_samples/perro/} y \texttt{audio\_samples/gallo/} respectivamente. Las tres grabaciones provienen de Wikimedia Commons bajo licencia libre (GFDL / CC BY-SA); la procedencia, autoría y licencia exacta de cada archivo se documenta en \texttt{audio\_samples/SOURCES.md}.

\begin{itemize}
\item \texttt{gato.ogg} — $44\,100$ Hz, $1.15$ s
\item \texttt{perro.ogg} — $44\,100$ Hz, $2.61$ s
\item \texttt{gallo.ogg} — $48\,000$ Hz, $3.84$ s
\end{itemize}

Nótese que estas tres grabaciones son considerablemente más cortas que los $10$ s asumidos en la Introducción para el conjunto original; gato y perro son incluso más cortas que la ventana de análisis de $3$ s, por lo que \texttt{analysis\_window} las procesa completas en lugar de recortarlas.

% ============================================================
% FIGURAS ANIMALES
% ============================================================

\begin{figure}[H]
\centering
\includegraphics[width=\linewidth]{gato.png}
\caption{Señal temporal (primeros 3 s) y espectro frecuencial correspondiente al sonido de un gato.}
\label{fig:gato}
\end{figure}

\begin{figure}[H]
\centering
\includegraphics[width=\linewidth]{perro.png}
\caption{Señal temporal (primeros 3 s) y espectro frecuencial correspondiente al sonido de un perro.}
\label{fig:perro}
\end{figure}

\begin{figure}[H]
\centering
\includegraphics[width=\linewidth]{gallo.png}
\caption{Señal temporal (primeros 3 s) y espectro frecuencial correspondiente al sonido de un gallo.}
\label{fig:gallo}
\end{figure}

\begin{figure}[H]
\centering
\includegraphics[width=\linewidth]{animales_comparacion.png}
\caption{Comparación de los espectros frecuenciales de los sonidos animales.}
\label{fig:animales_comparacion}
\end{figure}

% ============================================================
% INSTRUMENTOS
% ============================================================

\subsection{Instrumentos musicales}

Para estudiar diferentes registros frecuenciales se seleccionaron tres instrumentos musicales, con grabaciones reales descargadas de Wikimedia Commons:

\begin{itemize}
\item \texttt{contrabajo.oga} — contrabajo de jazz (\emph{walking bass}), $44\,100$ Hz, $17.74$ s; el instrumento entra recién a los $\sim\!2.5$ s, por lo que la mayor parte de la ventana de $3$ s analizada es silencio de estudio.
\item \texttt{piano.ogg} — piano, $44\,100$ Hz, $6.18$ s; una frase corta de varias notas, no una nota sostenida única.
\item \texttt{flauta.ogg} — primer movimiento lento de una sonata de Bach para flauta, $44\,100$ Hz (solo se procesan sus primeros $3$ s: introducción de clavecín y entrada de la flauta $\sim\!1$ s).
\end{itemize}

Ninguna de las tres es una nota aislada y sostenida como asumiría un análisis espectral idealizado; esto se retoma en la Discusión al interpretar sus frecuencias dominantes.

\begin{figure}[H]
\centering
\includegraphics[width=\linewidth]{contrabajo.png}
\caption{Señal temporal (primeros 3 s) y espectro frecuencial correspondiente al contrabajo. El instrumento entra a los $\sim\!2.5$ s.}
\label{fig:contrabajo}
\end{figure}

\begin{figure}[H]
\centering
\includegraphics[width=\linewidth]{piano.png}
\caption{Señal temporal (primeros 3 s) y espectro frecuencial correspondiente al piano.}
\label{fig:piano}
\end{figure}

\begin{figure}[H]
\centering
\includegraphics[width=\linewidth]{flauta.png}
\caption{Señal temporal (primeros 3 s) y espectro frecuencial correspondiente a la flauta.}
\label{fig:flauta}
\end{figure}

\begin{figure}[H]
\centering
\includegraphics[width=\linewidth]{instrumentos_comparacion.png}
\caption{Comparación de los espectros frecuenciales de los instrumentos musicales.}
\label{fig:instrumentos_comparacion}
\end{figure}

% ============================================================
% VOCES
% ============================================================

\subsection{Voces humanas}

Finalmente, se procesaron dos señales de voz sintetizada, generadas localmente con el motor Windows SAPI (\texttt{System.Speech.Synthesis}) usando dos voces instaladas distintas (\texttt{Microsoft David Desktop}, masculina, y \texttt{Microsoft Zira Desktop}, femenina), como sustituto reproducible mientras no se dispone de grabaciones humanas reales. \textbf{No son grabaciones de personas}; se etiquetan ``Persona 1'' y ``Persona 2'' únicamente para mantener la estructura de comparación de pares del resto del informe. Ambas voces leyeron el mismo texto, con el propósito de controlar el contenido lingüístico de las señales; al ser un motor en inglés, el texto en español se pronuncia con fonética inglesa, lo cual no afecta el análisis espectral pero sí impide comparar este resultado con locución en español real.

El texto sintetizado fue:

\begin{quote}
``Hoy estamos analizando senales de voz mediante la transformada de Fourier.''
\end{quote}

\begin{figure}[H]
\centering
\includegraphics[width=\linewidth]{persona1.png}
\caption{Señal temporal (primeros 3 s) y espectro frecuencial correspondiente a la Persona 1 (voz sintetizada, \texttt{Microsoft David Desktop}).}
\label{fig:persona1}
\end{figure}

\begin{figure}[H]
\centering
\includegraphics[width=\linewidth]{persona2.png}
\caption{Señal temporal (primeros 3 s) y espectro frecuencial correspondiente a la Persona 2 (voz sintetizada, \texttt{Microsoft Zira Desktop}).}
\label{fig:persona2}
\end{figure}

\begin{figure}[H]
\centering
\includegraphics[width=\linewidth]{personas_comparacion.png}
\caption{Comparación de los espectros frecuenciales de las dos voces humanas.}
\label{fig:personas_comparacion}
\end{figure}
